\cleardoublepage

\newrefsection

\chapter{外文翻译}

\begin{center}
    \zihao{-2} 端到端压缩网格单元渲染
\end{center}

\vspace{2em}

\sectionnonum{摘要}

在本文中，我们研究了使用现代 GPU 技术（尤其是网格着色器）对端到端压缩的三角网格进行渲染的方法。我们的方案能够在 GPU 内存中以压缩形式存储非结构化三角网格，并在着色器代码中即时解压，以便进行光栅化。

以往的典型方法通常仅在持久化存储和流式传输阶段使用压缩网格格式，但在提交渲染之前必须先将其解压到 GPU 内存中。相比之下，我们的方法在存储和 GPU 内存中均使用相同的压缩格式。因此，该压缩方法有效降低了存储超大三角网格所需的 GPU 内存，并且避免了由于 CPU 端解压缩导致的几何数据流式加载等待时间。端到端压缩还意味着，可以在 GPU 内存中完全驻留比以往更高几何细节的场景。

我们的方法基于对网格进行新颖的网格单元（Meshlet）分解，即将网格拆分为不相交的原始图元组，并分别进行压缩。利用网格着色器进行解压缩，使得在图元级别实现了事实上的随机访问，这对于选择性流式加载（Selective Streaming）和细粒度可见性计算（Fine-Grained Visibility Computation）等应用至关重要。

我们将该方法与多种常见的压缩网格单元格式在内存占用和渲染时间方面进行了对比。结果表明，我们的方法减少了 CPU-GPU 之间的内存带宽需求，而这正是外存渲染（Out-of-Core Rendering）中常见的瓶颈。

\section{引言}

开放世界游戏、科学模拟以及许多其他图形应用通常涉及巨大的场景，这些场景在存储、传输和渲染方面都面临挑战。如果场景无法完全容纳于 GPU 内存中，应用程序需要调度数据流式加载和解压缩，使数据能够恰好在需要时可用。对于这样的外部渲染（out-of-core rendering），最佳压缩率需要在内存占用和 GPU 友好性之间权衡。

传统上，渲染系统期望场景完全驻留（in-core）于 GPU 内存。如果场景大于 GPU 内存，则采用稀疏表示（sparse representation）并进行顺序访问（sequential access），可以实现外部渲染，即在内存中加载场景的空间局部区域（如果场景经过压缩，还需要进行解压）。此类顺序访问模式可用于分割、平滑或简化等任务，但通常不适用于渲染——绘制调用（draw calls）需要直接访问几何图元。如果场景由多个网格（mesh）构成，每个网格代表一个独立对象，则可以直接选择加载哪些网格。然而，若场景由一个超大网格或几个大型网格构成，使其无法完全驻留于 GPU 内存中，则通常的做法是将网格划分为多个区域。

这些区域应具备有限的空间范围，同时其图元数量应尽可能符合 GPU 的最佳处理方式。理想情况下，每个区域应对应于对象的一个有意义的部分 [CG08, CKLL09]，并且能够单独加载和解压缩。

流式渲染（streaming rendering）的需求早已被提出 [TMJ98, LH04]，但在当代游戏技术中仍然至关重要。微软的 DirectStorage API [Yeu20] 就是一个典型案例，它绕过 CPU，直接从非易失性存储（NVM）流式加载数据到 GPU 内存，并且最近扩展了硬件辅助 Lempel-Ziv 压缩解压 [Hoe22, CH22, Ura22]。Epic 的 Nanite 渲染引擎 [Epi20] 也采用了类似的方法。

然而，前述的外部渲染工具并未直接解决 GPU 内存占用问题。它们默认渲染基于未压缩表示，即着色器（shader）输入可以直接从内存读取。然而，鉴于 GPU 内存容量和带宽都是稀缺资源，我们更倾向于端到端（end-to-end）压缩渲染，即在着色器中直接解压缩三角形数据，使 GPU 内存仅存储压缩几何数据，而非未压缩数据。这样做的好处有：避免在 CPU 端进行单独的解压缩阶段；大幅减少几何数据的总体内存占用；简化引擎代码，无需同时维护压缩和未压缩的网格表示。

在本论文中，我们探讨了一种端到端压缩渲染架构，该架构直接在网格着色器（mesh shader） [Kub18] 中执行解压缩操作，处理称为网格单元（meshlet）的几何图元分组。能够在网格着色器中实时解压缩几何数据，使得构建流式渲染系统变得更加容易，因为 GPU 内存和非易失性存储中可以保持相同的压缩表示。

为了支持网格着色器中的解压缩，我们提出了一种适用于完全并行处理的拓扑连接表示（connectivity representation）。总结而言，我们的贡献包括：

\begin{itemize}

    \item 提出首个完整的端到端压缩网格渲染管线。
    \item 引入一种新颖的网格分解方法，将网格拆分为不相交的图元子集（称为laced wires），支持无冗余顶点的随机访问。
    \item 描述如何创建 laced wires，以及如何在 GPU 上高效地并行解码它们。
    \item 比较我们的方法与多种压缩和未压缩网格表示，分析其内存占用和渲染性能。

\end{itemize}

\section{相关工作}

本文关注三角网格（triangle meshes），它是计算机图形学中最广泛使用的网格表示方式。最常见的数据结构是索引三角网格（indexed triangle mesh），由一个存储每个三角形的三个连续顶点索引的三角形列表（triangle list）和一个存储顶点数据（例如位置或法线）的顶点列表（vertex list）组成。尽管这种数据结构可以直接被图形管线处理，但它并不紧凑。因此，已有大量研究致力于通过压缩顶点数据、连接信息，或两者同时压缩来减少存储开销 [MLDH15]。

大多数单速率（single-rate）网格压缩方法结合了三角形连接信息的无损编码与顶点数据的有损编码。为了简洁起见，本文不考虑有损连接信息压缩（即网格简化或重新网格化（re-meshing）），而是假设几何质量已经达到期望的压缩率。

\subsection{网格压缩（Mesh compression）}

顶点与属性压缩（Vertex and attribute compression）

在许多应用中，对顶点数据进行有损压缩是可以接受的，因为顶点坐标及其他属性通常不需要完整的浮点精度 [Dee95, Cal01, MSGS11, HV01]。 通过将坐标转换为定点格式并去除最低有效位来进行量化（quantization），通常不会导致可见的质量下降 [PBCK05, KXW*18]。事实上，某些网格（例如3D 扫描获取的网格）在其最低有效位中存储的往往只是噪声。然而，手工建模的网格可能包含不同部分的精度不一致，因此需要非均匀量化机制（non-uniform quantization）。

一种直接的解决方案是将网格进行分层分解，使每个部分内部具有均匀的精度。此外，还可以利用值序列的局部一致性（local coherence），通过增量编码（delta-encoding）或预测编码（prediction）进一步提高压缩率 [IA02]。

与位置量化（position quantization）不同，其他常见属性（如法线（normals） [CDE*14] 和纹理坐标（texture coordinates））的定义域和精度要求是已知的，因此它们的量化处理相对较少挑战。

连接信息压缩（Connectivity compression）

最紧凑的网格连接表示通常采用约束图遍历（constrained graph traversal），并使用低熵（low entropy）符号集合对遍历步骤进行编码。例如，三角形条带（triangle strips）是一种常见的方式 [Ise00]，它将顶点按序存储，每个新顶点与序列中的两个前序顶点共同定义一个新三角形。

标准的三角形条带模式中，每个新顶点会交替连接到前一个三角形的不同边，而在广义三角形条带（generalized triangle strips）中，使用特殊符号来选择要连接的边，从而构造更长的条带。进一步优化的遍历模式，如相邻条带排列（aligning strips） [Cho97] 或同心圆排列（concentric circles） [BPZ99]，可以进一步提高编码效率。

其他遍历策略包括生成树（spanning trees） [Tau98] 或区域扩展（region growing） [Ros99]，其中区域扩展可能是最具影响力的连接信息压缩方法。在区域扩展中，算法从网格的部分边界开始，以自定义的遍历策略逐步扩展已编码部分，并对新遇到的三角形进行编码。值得注意的遍历策略包括Cut-Border 机器（cut-border machine） [GS98]、EdgeBreaker [Ros99] 及其众多改进版本 [CR04, SKR01]。此外，区域扩展还可以通过顶点的度（valence）来表示，例如FreeLence [KPRW05]。

Laced Ring (LR) 结构 [GLRL11] 可以看作广义三角形条带的一种对偶表示（dual representation）。一个极长的三角形条带仅存储每个三角形所需的一个顶点引用，但在实际应用中，由于额外的操作（边交换、重启等），广义三角形条带会带来 30-40\% 的额外开销 [GLRL11]。Laced Ring 结构使用相邻顶点序列代替三角形序列。对于每条边，存储其两侧的 lace 引用，即相邻顶点组成的三角形。这种表示方式的存储开销仍然约为每个三角形一个引用，但相比三角形序列而言，更容易找到长的顶点序列（因为一个三角形平均只有3个相邻三角形，而一个顶点平均有6个相邻顶点 [GLLR13]）。本文将在此基础上进一步扩展 lace 的概念。

共享边界编码（Encoding shared boundaries）

如果一个大网格在压缩之前需要分割成多个区域，就必须处理不同区域间的共享边界，以避免裂缝（cracks）。最常见的避免裂缝的方法是沿着边界拆分网格，使不同区域间的顶点共享。但在这种情况下，共享的顶点需要在多个区域中被重复存储，这增加了内存占用，并可能增加保持顶点一致性的成本。

为了避免顶点重复存储，可以选择仅在一个区域中存储顶点数据，并在其他区域中通过引用访问 [YL07, CG08]，但这需要引用计数（reference counting）和独立的顶点加载（separate vertex loading）。

另一种方法是将边界顶点单独存储为顶点序列（wires）。Wires 可以基于现有的网格分割（subdivision）创建 [CKLL09]，或者通过递归网格分割的方式生成 [CH09]，即首先创建 wires，然后基于 wires 进一步划分区域。在后者的情况下，我们可以遍历分层网格分割结构，以提取任意大小的区域。

\subsection{GPU上的压缩域渲染（Compressed domain rendering on the GPU）}

在 GPU 上进行即时解压缩（on-the-fly decompression）的思想已经在一些涉及大规模数据集的领域中得到探索。这些领域的典型例子包括体绘制（volume rendering）[SW03, BCF03, FM07]，大规模地形数据的渲染 [DSW09, LC10]，以及空间加速结构 [LH07]。所有这些领域都涉及特殊的输入格式，其中的样本通常被组织成格点（lattices）、网格（grids）或体素块（boxes）。

相比之下，三角网格的连通性通常被表示为一个通用图（general graph），在使用单指令多数据（SIMD）并行计算时，需要更复杂的解压缩方法。典型的低熵（low-entropy）编码方案存在数据依赖性，每一步的解码都依赖于前一步的结果，因此并不适合 SIMD 计算。相反，我们需要一种半规则（semi-regular）的数据布局，使得多个线程可以直接访问压缩数据。这种布局可以采用长条带（long strips）[JBG17]，或者区域增长（region-growing）方案中的边界 [MSGS11]。

然而，在这些方法中，数据并行解压缩需要对 GPU 内存进行分散写入（scattered writes），这导致渲染必须经过一个额外的处理阶段，从而消耗大量的内存带宽。此外，网格的整体大小仍然受限于 GPU 内存的容量。商业游戏引擎（如 Unreal Nanite [KSW21]）在压缩率和解压缩速度之间进行了权衡，选择在 GPU 内存中解压缩到一种中间的、低熵但可随机访问的格式。

Nikolaev 等人 [NFR22] 在最近的一篇论文中探讨了与我们相似的想法，并比较了使用几何着色器（geometry shaders）和网格着色器（mesh shaders）进行即时解码和渲染的方法。然而，与我们不同的是，他们采用了一种区域增长（region-growing）方案 [Ros99]，这在 GPU 上存在严重的性能缺陷：不仅区域增长方式破坏了数据引用的局部性（locality of data references），其解码方案（基于网格连通性的区域增长和几何数据的增量编码）本质上也是串行的。因此，该方法只能在网格单元（meshlet）级别支持并行解码，而无法在单个三角形级别上并行解码。不幸的是，这一重大限制并未在该论文中得到深入讨论或评估。我们将在第 5 章展示我们的方法相比 Nikolaev 等人方法的性能提升。

\subsection{网格布局优化}

在标准的深度缓冲（depth buffer）硬件渲染中，图元（primitives）可以以任意顺序提交到图形管线，因此我们可以自由选择最适合特定需求的网格布局。引用的局部性（locality of references）通常是优化的主要关注点，因为它带来了多个好处。首先，它可以提高 GPU 顶点缓存（vertex cache）的访问一致性，从而减少冗余的逐顶点计算 [Dee95, Hop99]。许多算法被提出用于生成针对缓存优化（cache-aware）[LY06, CK07] 或缓存无关（cache-oblivious）[YLPM05] 的网格布局。缓存优化算法需要准确的缓存大小信息，而缓存无关算法则在不了解缓存具体细节的情况下优化局部性。例如，可以对三角形到顶点的引用矩阵进行排序，以最小化带宽消耗 [IL05]。即使是简单的技术，比如遵循空间填充曲线（space-filling curve），也能生成接近理论最优的网格布局 [VSSP12]。

除了引用的局部性之外，其他标准也可以在布局生成时考虑，例如紧凑的内存占用、几何一致性（geometric coherence）或最小化过度绘制（overdraw）[CSN*12, SNB07, HS16]。

网格布局也会在网格被编码时，受到压缩算法的隐式影响。深度优先遍历（depth-first traversal）可以产生长序列的相连三角形用于压缩，如果采用“缠绕”（winding）路径，则可能与缓存无关算法相似 [CK07]。然而，非缠绕的序列可能具有同样的压缩率 [GLLR13]，但其缓存一致性可能较差 [DBGP05]。同样，采用区域增长（region growing）[Ros99] 的广度优先遍历（breadth-first traversal）会引用整个网格，而不一定具有良好的缓存一致性。换句话说，良好的压缩率并不意味着良好的缓存一致性。

此外，现代 GPU 以批处理（batches）的方式工作，即一组包含一定数量顶点和三角形的图元组。要实现最佳吞吐量，不仅需要考虑数据引用的位置，还需要一个符合批处理需求的布局 [YL06]。可以通过贪心拆分（greedy splitting）缓存无关布局来计算批次，例如在运行时最小化过度绘制 [SNB07]。然而，如果网格在预处理时已经被分割，例如为了生成支持随机访问的压缩网格表示，我们希望选择的网格分段能够同时作为 GPU 友好的批次，以便提升渲染效率。

\section{压缩网格单元表示}

我们提出 Laced Wires，这是一种可随机访问的压缩网格结构，结合了先前方法的优点，即 Laces [GLRL11] 和 Wires [CKLL09, CH09]。我们的表示既紧凑又适用于 SIMD，并且能够在高 SIMD 利用率下实现端到端压缩，这是至关重要的。

Wire 本质上是由边连接的顶点序列 [CKLL09]。我们利用 Wires 来表示网格单元的边界和内部，使得相邻的网格单元在共享边界时可以引用相同的 Wire。而 Laces 则存储了某条边左侧和右侧两个三角形的顶点信息。我们将 Laces 应用于 Wires，使相邻的 Laced Wires 相互交错，而不会产生冗余顶点。

Laced Ring（LR） 结构使用尽可能长且曲折的环状顶点序列。LR 方案的优化目标是提升 CPU 内存中的全局存储效率，但因此牺牲了局部引用的特性。而我们的目标是在 GPU 的 SIMD 单元缓存中获取和解码有限数量的内存，以完整描述一个网格单元。因此，我们的表示采用 每个网格单元由多个短序列组成，而不是像 LR 那样使用单个长序列。相邻 Laces 的交错设计避免了顶点的重复存储，使得该格式没有冗余数据，并且可以防止三角形之间出现裂缝。图 1 展示了 LR 与 Laced Wires 之间的对比。

\begin{figure}[h]
    \centering
    \begin{minipage}{\linewidth}
        \centering
        \includegraphics[width=\linewidth]{1.png}
    \end{minipage}

    \vspace{0.5em}
    
    \begin{minipage}{\linewidth}
        \fontsize{11pt}{13.5pt}\selectfont
        \bfseries
        图1\quad（左图）Laced Ring（LR）由单一的闭合顶点序列（黑色）组成。附着在 LR 序列边上的三角形以蓝色和黄色显示，仅用于更好地区分。部分孤立三角形（灰色轮廓）仍然存在；其中大多数三角形可以通过特殊编码集成，而无需额外的存储开销。（右图）相同的网格被分解为三个网格单元（左上、右上、底部），其边界由六条外部 wire 标示（彩色折线：蓝色、橙色、绿色、红色、黄色、洋红色）。三个内部 wire（黑色折线）使用类似螺旋的哈密顿路径贪心构造。与 wire 关联的相邻三角形使用与该 wire 相同的颜色显示。角点顶点以黑点表示。
    \end{minipage}
\end{figure}

我们压缩表示的构建包括四个步骤：网格单元构建（第 3.1 节） 将三角形分配到 网格单元；连通性编码（第 3.2 节） 将网格单元的三角形数据转换为 laced wire 格式；顶点编码（第 3.3 节） 在局部坐标框架中，考虑用户定义的误差容限，确定每个顶点的最小比特率；数据打包（第 3.4 节） 将数据整理成二进制格式，以完成构建过程。

\subsection{网格单元构建}

我们将三角形聚类为互不重叠的网格单元，以确保不超过原始几何体的极限：每个网格单元最多包含 84 个三角形和 64 个顶点，按照 Kubisch [Kub18] 的推荐。理想情况下，网格单元是一个连通的、紧凑的、平坦的子网格。我们预期，如果一个网格单元是平坦的，其几何体的“垂直”维度可以更好地压缩，而且形状紧凑的网格单元将整体表现出更好的引用局部性。

我们的网格单元构建算法基于使用四面体误差度量 [GWH01] 的分层面聚类。起初，所有的三角形都是单独的聚类，即每个聚类只包含一个三角形，之后聚类会通过迭代方式合并。该方法基于网格的对偶表示，聚类作为节点，节点之间通过加权边连接，表示相邻的聚类。只有共享至少一条边的聚类才会被认为是相邻的。对偶边的折叠对应着两个相邻聚类的合并。在每次合并后，我们删除所有与新聚类相邻的对偶边，这些边在折叠后会导致聚类超过最大顶点数或三角形数。我们在无法继续对任意两个相邻聚类进行合法合并操作时终止网格单元生成，即当对偶网格中没有更多对偶边时。

\subsection{连接性编码}

为了编码网格单元的连接性信息，我们引入了 laced wires。我们采用了“laces”[GLRL11] 的思想，即指定与边序列相邻的三角形，而不是三角形序列。我们将 laces 应用于两个相邻网格单元共享的顶点的边，并将这种结构称为外部 wire。每个网格单元的内部也被转换为一个 laced wire；得到的结构称为内部 wire。
每个网格单元与 k 个邻居相关，由一个内部 wire 和 k 个外部 wire 组成。在图 1 左侧，展示了 LR 的一个示例，它由一个单一的顶点循环组成。图右侧展示了我们的数据结构：网格被分解为三个网格单元，导致六个外部 wire（彩色）和三个内部 wire（灰色）。

与网格分解为网格单元的方式不同，后者沿边界具有重复的顶点，laced wire 表示法避免了冗余。每个三角形只存储在一个 wire 中；每个顶点只存储在一个 wire 中，除了三个或更多网格单元共享的角落顶点。然而，网格单元的几何体可以完全从与该网格单元 相关的 wires 中重构，从而实现随机访问。

wire 中顶点的顺序遵循边的序列。在外部 wire 中，边序列的方向沿着相邻网格单元 的边界按逆时针方向设置，且相邻网格单元的索引较小。对于内部 wire，使用一种贪心算法确定 Hamiltonian 路径，该路径类似于一个侵蚀过程（图 1 右侧）。如果找不到 Hamiltonian 路径，则通过选择多个尽可能长的路径来近似路径。为了找到 Hamiltonian 路径，顶点根据其到边界的最短路径排序。路径的构建通过迭代地添加一条新边，将一个顶点连接到路径的当前末端，并始终优先选择较低级别的顶点。这一策略以类似螺旋的模式从边界征服网格单元内部。如果没有剩余的自由顶点可以连接到路径的末端，可以通过中介“虚拟”边连接任何自由顶点，并插入一个特殊符号，表示序列重启。

外部 wire 的 laces 引用相邻内部 wire 中的顶点，反之亦然。外部 wire 使用相邻内部 wire 的本地编号来存储对内部 wire 中顶点的引用。相反，内部 wire 必须引用多个外部 wire。为了实现随机访问，每个对外部 wire 顶点的引用必须能够在不遍历整个边界的情况下解析。因此，我们定义了一个边界顶点的唯一排序，从与邻近网格单元共享的角落开始，按逆时针顺序进行排序。对外部 wire 的引用使用此排序，并通过一个并行索引重写步骤解析，基于沿边界的角落序列（因此也包括外部 wire）。

在最佳路径中，每个三角形与路径的一个边相邻，可以编码为一个 1 引用三角形（1RT）。不幸的是，最佳路径很难找到，因此我们必须处理两种额外的配置 [GLRL11]，如图 2 所示。首先，与两条边相邻的二引用三角形（2RT）必须避免重复编码，因此我们需要使其中一个引用无效。其次，与任何边都不相邻的三角形不能省略，因此必须编码为三引用三角形（3RT）。所有 3RT 都分配给内部 wire；即，外部 wire 中不能存在不规则三角形。为了减少引用数量，具有共享非路径边的 3RT 和 2RT 可以通过 2RT 的两个引用一起编码。

\begin{figure}[h]
    \centering
    \begin{minipage}{\linewidth}
        \centering
        \includegraphics[width=\linewidth]{2.png}
    \end{minipage}

    \vspace{0.5em}
    
    \begin{minipage}{\linewidth}
        \fontsize{11pt}{13.5pt}\selectfont
        \bfseries
        图2\quad 解码过程中区分的三角形配置。从左到右：需要一个引用的常规三角形、需要两个引用的角落三角形、需要三个引用的未连接三角形、相邻的二/三引用三角形对。
    \end{minipage}
\end{figure}

\subsection{几何编码}

我们通过在每个 wire 的局部坐标系中均匀量化每个顶点来进行编码。局部坐标系直接决定了量化域。因此，我们试图识别一个量化域，以最小化每个顶点所需的位数，如下所述。

Per-wire 局部坐标系（Per-wire local coordinate frame）

首先，我们为每个 wire 构建坐标系。由于我们的网格集是在考虑平面的基础上构建的，因此我们可以通过将其与网格集的主平面对齐来找到一个紧密拟合的有向包围盒（OBB）。OBB 规范进一步量化为九个值（每个维度的旋转、平移和缩放），每个值占用一个字节。

顶点量化（Vertex quantization）

我们使用 wire 的 OBB 来计算顶点位置的实际量化。我们的量化方案允许在单个 wire 的粒度上选择所需的精度，这对于忠实表示具有细小细节的网格，或支持每个对象多个细节级别表示可能很重要。我们确保从量化的顶点中可以恢复完整的顶点精度（详细信息请参见补充部分）。corner，即由多个网格集共享的顶点，会在整个网格的坐标系中进行编码，并作为伪线段（pseudo-wire）单独存储（没有连接性）。请注意，每个顶点仅存储一次，即使经过重度量化，也不可能出现裂缝。

\subsection{二进制格式}

每个laced wire都表示为一个完全自包含的位串，包含解码所需的所有信息，且支持随机访问解码。我们将数据区分为头部数据、连接数据和顶点数据（图3）。头部包含顶点数量和 OBB。连接数据将左右侧的编织数据存储在单独的子序列中，因为只需要其中一个就可以解码网格集的外部 wire。每个符号大小为一个字节，其中两个比特用于编码三角形类型，六个比特用于区分最多64个顶点。顶点数据（量化后的 x、y、z 值）存储在打包的位串中。

\begin{figure}[h]
    \centering
    \begin{minipage}{\linewidth}
        \centering
        \includegraphics[width=\linewidth]{3.png}
    \end{minipage}

    \vspace{0.5em}
    
    \begin{minipage}{\linewidth}
        \fontsize{11pt}{13.5pt}\selectfont
        \bfseries
        图3\quad laced wire的二进制表示。头部（蓝色）包含顶点数量 𝑁、偏移量 𝑂指向二进制表示中的顶点以及局部 wire 变换𝑇。连接数据（绿色）被拆分为与左侧、右侧和不规则连接对应的部分，每部分前都有相应的大小。顶点数据（灰色）包含每个维度的15个位用于量化位，随后是量化后的顶点数据。
    \end{minipage}
\end{figure}

除了编码的 wire，我们还需要元数据来查找每个网格单元的 wire：我们需要其内部wire、外部 wire 和 corner 的索引。我们将这些信息存储在每个网格单元的结构中，并包括指向该网格单元第一个外部 wire 索引的偏移量。

\section{GPU 解码器实现}

我们根据需要从非易失性存储中预取压缩的网格块数据。与传统的几何体预加载不同，预取的数据以压缩形式存储在 GPU 内存中，从而显著减少了内存占用。通过结合以下两种方式带来的内存效率提升：(1) 场景数据的稀疏驻留，(2) 内存中的压缩表示，可以处理单独使用任一方法时无法处理的巨大场景。

在解码过程中，我们使用一个包含放大着色器和网格着色器的管线，如图 4 所示。前者用于决定哪些网格块需要在随后的网格着色器中解压缩以进行渲染。网格着色器允许我们完全自由地重新定义 GPU 管线的几何阶段。我们展示了如何利用网格着色器进行压缩渲染，使得几何阶段的输入和输出都能够节省内存带宽。输入带宽通过仅加载压缩网格块来得到保留。最终的三角形、顶点和属性解压缩是在网格着色器中实时完成的。因此，未压缩的网格流的完整带宽仅在网格着色器和光栅化阶段之间需要，且该带宽由硬件缓存提供支持。

\begin{figure}[h]
    \centering
    \begin{minipage}{\linewidth}
        \centering
        \includegraphics[width=\linewidth]{4.png}
    \end{minipage}

    \vspace{0.5em}
    
    \begin{minipage}{\linewidth}
        \fontsize{11pt}{13.5pt}\selectfont
        \bfseries
        图4\quad 在这个概述中，颜色表示涉及的内部 wire 数据（绿色）、外部 wire 数据（蓝色）或两者都有。每个放大着色器实例通过四个线程的组启动，每个线程决定一个网格块的可见性。根据可见性，启动 0 到 4 个网格着色器实例。每个网格着色器实例使用 32 个线程执行，从其内部线和外部线解码单个网格块。
    \end{minipage}
\end{figure}

\subsection{放大着色器}

放大着色器用于选择当前帧所需渲染的网格。基于集群的剔除技术是众所周知的，可以让我们在管线中更早地丢弃更多的几何体 [KMGL99, HA15, Wih16]。第一步应用视锥体剔除。我们使用内部 wire 的局部坐标变换来定义网格块的包围球。我们可以使用包围盒的中心作为包围球的中心，包围盒对角线的一半作为包围球半径的初步估计。由于包围体会被频繁测试，我们发现包围球比包围盒更快，因为它只需要转换一个中心点，而不是八个角点。我们使用每个网格块 8 位来存储初始半径的比例因子，以确保缩放后的内部 wire 包围球包含整个网格块。为了执行视锥体剔除，我们将中心和缩放后的半径从局部坐标转换到全局坐标系，并将其与六个视锥体平面进行测试。

第二步应用背面剔除，使用一个覆盖网格块所有法线的圆锥。由于网格块在构建时已经非常平坦，大多数法线会大致指向局部坐标系的 x 轴。因此，我们不显式存储圆锥轴，而是使用内部 wire 的局部坐标系的 x 轴。我们使用内部 wire 包围盒的中心作为圆锥的顶点，与之前包围球的做法类似，并使用每个网格块 8 位来存储圆锥角度，确保它包含网格块的所有三角形法线。角度超过 90º意味着网格块非常弯曲，应该将其拆分为多个网格块。为了确保圆锥顶点位于所有三角形的负半空间，我们使用另 8 位存储圆锥顶点相对于中心沿圆锥轴的偏移量（图 5）。

\begin{figure}[h]
    \centering
    \begin{minipage}{\linewidth}
        \centering
        \includegraphics[width=\linewidth]{5.png}
    \end{minipage}

    \vspace{0.5em}
    
    \begin{minipage}{\linewidth}
        \fontsize{11pt}{13.5pt}\selectfont
        \bfseries
        图5\quad 在这个概述中，颜色表示涉及的内部 wire 数据（绿色）、外部 wire 数据（蓝色）或两者都有。每个放大着色器实例通过四个线程的组启动，每个线程决定一个网格块的可见性。根据可见性，启动 0 到 4 个网格着色器实例。每个网格着色器实例使用 32 个线程执行，从其内部 wire 和外部 wire 解码单个网格单元。
    \end{minipage}
\end{figure}

通过内部 wire 的局部变换，我们可以执行视锥体剔除和背面剔除，每个网格体只需要额外的 3 字节数据，用于存储包围球半径的缩放因子、圆锥角度和顶点偏移。如果网格体没有被剔除，则启动一个网格着色器实例进行解压。

\subsection{网格着色器}

每个网格着色器实例都通过一个线程组（32 个线程）启动，负责解压单个网格体。内部和外部 wire 的顶点需要进行反量化，并从 wire 的局部坐标系转换到全局坐标系。此外，三角形索引需要从内部 wire 和每个相关外部 wire 的一侧的 laced wire 表示中解码。网格着色器需要在写入输出数组之前指定输出顶点和三角形的数量。因此，我们遍历所有 wire（内部和外部），并将它们的头部加载到共享内存中。我们从头部提取出总的顶点数，并将其传递给网格着色器。不幸的是，网格体中的三角形数量不能以相同的方式确定，因为它不能直接从描述 wire 连接性的符号数量推断出来（其中包括不规则情况，如空引用和重启）。因此，我们总是将网格着色器输出的三角形数量设置为最大三角形数，即 84，并用零填充所有未使用的三角形。我们通过实验证明，零填充比显式存储和填写准确的三角形数量更快。

内部 wire 连接性

对于内部 wire，我们必须解码 wire 连接性的两侧以及不规则的三角形。我们为 wire 的每一侧分配一个线程。每个线程读取一个符号，并确定它是否编码了一个三角形。线程投票形成一个位掩码，指示哪些三角形存在。通过计算掩码中设置的位数，每个线程可以确定它在输出数组中的写入偏移量。遇到有效三角形的线程解码其顶点并将结果写入输出数组。每个线程然后读取一个新的符号，直到输入被消耗完。最后，不规则的 3RT 连接性由每个三角形一个线程解码。

内部 wire 顶点

内部 wire 顶点与每个顶点一个线程并行解码。首先，每个线程确定每个维度的量化位数，可以在顶点数据的前 15 位找到。根据顶点大小，每个线程计算它在编码 wire 顶点中的偏移量，并将量化值解包到寄存器中。从那里，顶点被反量化，得到线的局部坐标系中的位置。变换信息从存储在共享内存中的头部解包，顶点被转换到全局坐标系中。应用视图-投影变换后，顶点被写入输出。

外部 wire 连接性

解码外部wire的连接性类似于内部wire，但必须考虑多个外部wire，而不是单个内部wire。每个外部wire包含不同数量的符号，使得并行解码符号变得更加困难。我们不再像以前那样在单一wire的边缘上并行化，而是对多个wire进行并行化。外部wire的解码分为两个阶段。第一阶段，我们设置并行解码的工作。我们遍历外部wire并确定需要解码哪个方向的连接性。对于每个外部wire的每个符号，保留一个线程。第二阶段，符号被并行解码，直到没有符号可以解码为止。像以前一样，线程投票决定它们是否遇到有效的三角形，以形成一个位掩码。位掩码让线程将解码的三角形写入正确的位置。如果我们无法将单条wire的所有符号分配给线程，我们将对wire的连接性进行多次迭代解码。

外部 wire 顶点

解码外部wire顶点遵循与解码内部wire顶点相同的过程。量化值从其二进制表示中解包，反量化到局部坐标系中给定的位置，随后使用存储在wire头部中的变换转换到全局坐标系中。单个外部wire通常没有足够的顶点来利用网格着色器线程组的所有线程。因此，我们对具有不同数量顶点的多个wire进行并行化。多个wire的联合解码需要对每个wire的顶点数进行前缀和运算，以获得顶点输出数组的偏移量。在前缀和过程中，每个线程确定一个顶点的工作描述，即分配给该线程的外部wire的偏移量以及wire索引。然后，我们遍历各wire并分配线程到顶点，直到线程组完全分配。接下来，我们使用寄存器洗牌将工作描述传递给解码线程。每个解码线程计算其wire的顶点大小，找到它的顶点偏移量，解码它并将其写入输出数组，位置由偏移量和顶点在该wire中的索引决定。如果wire位于网格体的右侧，则顶点顺序被反转，以跟随网格体边界。线程组继续解码下一个批次的顶点，直到所有wire都解码完毕。

两个相继的外部wire构成网格体的边界，它们共享一个角点。因此，我们可以为每条外部wire分配一个线程，用于检索起始角点并将其写入输出数组的相应位置。注意，只有当前处理的网格体的角点需要被加载。此外，所有角点都编码在相同的全局坐标系中，因此它们的变换可以存储在常量内存中。

\section{评估}

我们从空间（压缩率）和时间（解码性能）两个方面评估我们的方法，并将其与其他压缩和未压缩的网格格式进行比较。

\subsection{网格单元格式}

为了研究网格单元几何和连通性压缩的影响，我们将我们的 laced wire 格式与四种不同的表示形式进行比较。每种表示形式采用两种不同的内部格式（基础版和紧凑版），最终形成六种候选格式进行对比。所有基础版（basic）均使用全精度 32 位顶点坐标和 32 位图元索引。

三角形集合网格单元（TS）

最简单的格式之一是将每个网格单元表示为一个三角形集合（triangle soup），其中每个三角形都包含其所有顶点的副本。尽管这种格式的表示最不紧凑，但它不需要进行解码。在网格单元格式中，它作为一个基准，用于更好地理解数据大小和解码开销之间的权衡。要加载一个网格单元，只需要其在全局缓冲区中的偏移量以及图元数量。

值得注意的是，尽管 TS 格式很简单，但它仍然可以被视为一种网格单元格式，因为它局部地对图元进行了分组。我们没有比较全局三角形集合格式，因为我们认为它过于基础。相反，我们将其与全局索引网格（GI）进行比较（见下文）。

在该格式的紧凑版本中，顶点根据整体网格包围盒的大小进行全局量化。量化后的顶点可以被紧凑地存储在内存中，使每个顶点占用的比特数正好匹配其量化精度。同时，网格单元数据会填充到下一个 32 位字边界。

索引网格单元（IX）

IX 格式通过使用索引缓冲区，避免了 TS 格式中顶点的冗余复制。IX 依赖于全局顶点缓冲区来存储唯一的网格顶点。网格单元包含一个顶点索引缓冲区，该缓冲区存储了对全局顶点缓冲区的引用，以标识网格单元中的各个顶点。此外，每个网格单元还有一个图元索引缓冲区，该缓冲区定义了局部三角形，即通过引用局部顶点索引缓冲区的方式表示。通过这种从图元到局部顶点索引，再到全局顶点的两级间接引用，该格式提高了数据的紧凑性。

在该格式的紧凑版本中，我们同样对顶点进行全局量化，并将其紧凑存储。此外，每个图元对顶点索引缓冲区的引用仅占 8 位，因为网格单元最多包含 64 个顶点。顶点索引缓冲区的条目占 32 位，以确保能够寻址大规模网格的顶点。

独立网格单元（SC）

SC 格式通过引入一定的顶点冗余，使网格单元成为独立的、自包含的单元。与 IX 格式不同，SC 格式没有局部顶点索引缓冲区，而是每个网格单元直接存储其自身的顶点。边界上的顶点会被复制。为了访问其数据，每个 SC 记录存储从图元索引和顶点数据的起始偏移量，以及网格单元中三角形和顶点的数量。

在紧凑版本中，顶点被量化并紧凑存储，同时网格单元的三角形索引与 IX 格式相同，占 3 字节。

全局索引网格（GI）

除了不同的网格单元表示形式外，我们还添加了 GI 格式的数据进行对比，这是最广泛使用和最知名的网格格式。GI 由一个全局顶点缓冲区和一个全局索引缓冲区组成，后者通过引用顶点缓冲区来定义三角形的连接关系。

在连接性存储方面，GI 需要每个三角形占 12 字节。由于需要寻址较大的顶点缓冲区，因此无法对索引进行紧凑存储，也就无法实现压缩。然而，顶点可以以基础格式存储（每个顶点 12 字节），或者采用全局量化格式（见 5.2 节）。

\subsection{评估设置}

所有实验均在配备 64 GB 内存和 NVIDIA RTX 4090 的 Intel Core i7-7700 处理器上进行。我们使用了三个不同大小的场景，分别标记为 S、M 和 L，它们由各种不同的网格组成，如图 6 所示。场景 L 包含所有网格，而场景 M 和 S 是其子集，具体信息详见表 1。所有格式使用相同的网格分解方式，仅在网格单元的表示方式上有所不同。

\begin{figure}[h]
    \begin{minipage}{\linewidth}
        \fontsize{11pt}{13.5pt}\selectfont
        \bfseries
        表1\quad 表中 S、M 和 L 场景的列表示某个特定网格是否包含在相应场景中。此外，表中给出了网格中的三角形数量（$n_{t}$）、顶点数量（$n_{v}$）和网格单元数量（$n_{m}$），以及每个网格单元的平均三角形数量（$\overline{n_{t}}$）和平均顶点数量（$\overline{n_{v}}$）。
    \end{minipage}

    \vspace{0.5em}

    \centering
    \begin{minipage}{\linewidth}
        \centering
        \includegraphics[width=\linewidth]{6.png}
    \end{minipage}
\end{figure}

\begin{figure}[h]
    \centering
    \begin{minipage}{\linewidth}
        \centering
        \includegraphics[width=\linewidth]{6.jpg}
    \end{minipage}

    \vspace{0.5em}
    
    \begin{minipage}{\linewidth}
        \fontsize{11pt}{13.5pt}\selectfont
        \bfseries
        图6\quad 我们的测试数据包括由以下大型网格组成的三个场景。从左到右依次为：Armadillo、Dragonfly、Neptune、Dragon、Statuette、Crab（雌性蓝蟹，致谢史密森学会）、Bee（Eulaema Meriana 蜂，致谢史密森学会）、David（致谢数字米开朗基罗项目）。在每张图像的右下半部分，网格单元以随机颜色编码进行标示。
    \end{minipage}
\end{figure}

对于输入网格的顶点（在编码之前），我们通过实验证明了一种全局量化方法，该方法可以为所有备选格式提供视觉上令人满意的结果，并且我们的方法配置为低于由此产生的最大解码误差。大多数网格的顶点坐标均被量化为 14 位，但 Armadillo 采用 11 位，Dragonfly 采用 13 位，而 David 采用 16 位。

\subsection{压缩率与解码时间}

为了评估压缩率，我们将我们的方法所需的内存占用与其他备选格式（TS、IX、SC 和 GI）的基本版和紧凑版进行比较。表 2 展示了各方法的内存需求和每帧渲染时间。相较于基本版，紧凑版的压缩比例分别为 TS 的 2.23倍，IX 的 2.38倍，SC 的 2.91倍，GI 的 1.22倍。

除了 TS 以外，优化带来的计算开销几乎不会对每帧渲染时间产生显著影响（IX 1.13倍帧率，SC 0.99倍帧率，GI 1.00倍帧率）。然而，在 TS 格式中，所应用的优化几乎使所有场景的平均帧时间翻倍（导致帧率为原本的0.56倍）。

\begin{figure}[H]
    \begin{minipage}{\linewidth}
        \fontsize{11pt}{13.5pt}\selectfont
        \bfseries
        表2\quad 在不同场景大小下，使用基本版和紧凑版格式所需的空间（$s_{b}$、$s_{c}$）以及每帧渲染时间（$t_{b}$、$t_{c}$），并与 laced wires 方法进行比较。此外，表中还展示了紧凑版相较于基本版的压缩率（$s_{b}/s_{c}$）和相对渲染时间（$t_{b}/t_{c}$）。下标$b$表示基本版，下标$c$表示紧凑版。
    \end{minipage}

    \vspace{0.5em}

    \centering
    \begin{minipage}{\linewidth}
        \centering
        \includegraphics[width=\linewidth]{7.png}
    \end{minipage}
\end{figure}

我们将紧凑格式对 TS 方式帧时间的显著影响归因于以下事实：在该格式中，几乎所有的解码工作都涉及顶点处理，因为其连通性是隐式的，而需要解码的顶点数量远多于其他格式。因此，解包和去量化顶点引入了较大的解码开销。此外，由于顶点坐标通常不再对齐至 4 字节地址，导致内存传输量增加。值得注意的是，其他方法通过减少 GPU 内存传输来摊销解码工作，而 SC 甚至在性能上略胜一筹。

在表 2 中比较紧凑版本的绝对解码时间可以发现，在与 laced wires 竞争的方法中，SC 既实现了最佳的压缩率，也达到了最快的解码时间。SC 以 807 MB 的存储空间编码了大型场景中的 1.09 亿个三角形和 5400 万个顶点，并在 3.3 毫秒内渲染一帧。相比之下，GI 无论在基础版本还是紧凑版本中都需要大约 3 倍长的时间。

在压缩率方面，laced wires方法优于所有竞争格式，仅占下一个最佳压缩格式 SC 所需存储的 63\% 。然而，这种额外的压缩率是以较高的解码成本为代价的：解码和渲染整个大型场景需要 8.18 毫秒，大约是 SC 的 2.5 倍，但仍比标准格式 GI 快 28\%。

为了更清楚地了解我们方法中的内存需求分布情况，我们在表 3 和表 4 中展示了三种测试场景的详细数据。内部和外部 wire 对数据大小的贡献大致相等，总共占总存储需求的 80\% 左右。角点数据仅占大约 3–4\% 的存储空间，尽管所有顶点都在一个较大的全局域内进行量化，因此需要更多的比特来满足精度要求。邻接信息（即将外部 wire 与网格单元的内部 wire 关联，并将每个外部 wire 与两个 corner 关联的数据）大约占总数据量的 18\%。

\begin{figure}[h]
    \begin{minipage}{\linewidth}
        \fontsize{11pt}{13.5pt}\selectfont
        \bfseries
        表3\quad 我们方法中场景划分的详细信息。对于内部和外部 wire，我们列出了 wire 数量（$w_{i}$、$w_{e}$）、三角形数量（$t_{i}$、$t_{e}$）、顶点数量（$v_{i}$、$v_{e}$）以及角点数量（$c$）。下标$i$表示内部 wire，而$e$表示外部 wire。
    \end{minipage}

    \vspace{0.5em}

    \centering
    \begin{minipage}{\linewidth}
        \centering
        \includegraphics[width=\linewidth]{8.png}
    \end{minipage}
\end{figure}

\begin{figure}[h]
    \begin{minipage}{\linewidth}
        \fontsize{11pt}{13.5pt}\selectfont
        \bfseries
        表4\quad 内部 wire 数据（$d_{i}$）、外部 wire 数据（$d_{e}$）和 corner 数据（$d_{c}$），以及解码和渲染整个场景所需的总数据量（$d_{t}$）。数据中包含了对齐至 32 位边界的填充（通常会产生 1 - 4\% 的额外开销）。邻接数据（$d_{a}$）用于存储每条内部 wire 对应的外部 wire 标识符，以及每条外部 wire 对应的两个corner 索引。
    \end{minipage}

    \vspace{0.5em}

    \centering
    \begin{minipage}{\linewidth}
        \centering
        \includegraphics[width=\linewidth]{9.png}
    \end{minipage}
\end{figure}

\subsection{连接关系的压缩率}

我们对连接关系所需的内存进行单独分析，结果如表 5 所示。为了简洁起见，本分析仅针对大场景进行研究。

\begin{figure}[h]
    \begin{minipage}{\linewidth}
        \fontsize{11pt}{13.5pt}\selectfont
        \bfseries
        表5\quad 不同方法在原始格式（$s_{b}$）和紧凑格式（$s_{c}$）下的连接关系与几何数据的内存消耗，紧凑格式相较于基础格式的相对压缩率（$s_{b} / s_{c}$），以及紧凑格式相较于紧凑型全局索引网格（GI）的基准压缩率（$g_{c} / s_{c}$）。此外，我们还列出了基础格式和紧凑格式下每个元素（$e_{b}$、$e_{c}$）的内存占用，其中，连接关系的元素为三角形，而几何数据的元素为顶点。
    \end{minipage}

    \vspace{0.5em}

    \centering
    \begin{minipage}{\linewidth}
        \centering
        \includegraphics[width=\linewidth]{10.png}
    \end{minipage}
\end{figure}

TS 格式并不需要显式存储连接关系数据。IX 和 SC 的原始方法在使用图元索引定义连接关系方面与 GI 相似，但额外需要存储每个网格单元的三角形数量以及图元索引的偏移量。与 GI 不同的是，IX 和 SC 的图元索引仅需 8 位即可表示，因为它们仅引用局部顶点，这带来了 3.88 倍的压缩率。

对于 laced wires 方法，内存需求主要来自于编码 lace 的引用。我们的方法需要存储无法明确归类为几何数据或连接关系数据的信息，即 邻接数据。我们将这些数据相对于场景中的三角形数（$n_{t}$）和顶点数（$n_{v}$）进行拆分，并分别归类为连接关系和几何数据。其中，$n_{t}/(n_{t} + n_{v})$的通用数据被分配到连接关系部分。

为了建立统一的对比基准，我们以 GI 的连接关系数据作为参考，计算压缩率。可以看出，IX 和 SC 的压缩率为 3.84，即每个三角形的存储需求为 3.13 字节，而它们各自的基础版本为 12.13 字节，GI 则需要 12 字节。相比之下，我们的方法相较于 GI 取得了 5.96 的连接关系压缩率，相当于每个三角形仅需 2.01 字节。这一结果表明，laced wires 方法的连接关系比 IX 和 SC节省了 1.56 倍。

\subsection{几何数据的压缩率}

GI 的几何数据占用的内存比网格单元格式（除了 laced wires 之外）更少。这是因为 TS 在每个三角形内都会重复存储顶点，IX 需要为每个网格单元维护一个本地顶点索引数组，并在全局顶点数组中存储一个 4 字节索引，而 SC 在每个相邻的网格单元之间都会存储边界顶点的副本。

在 laced wires 方法中，几何数据的内存需求计算包括所有用于解码顶点的数据，例如 指向 wire 内第一个顶点字节的偏移量、顶点数量以及每个网格单元顶点在各维度上的量化位数。此外，每个 wire 还包含 9 字节的局部变换数据。

Laced wires 以1.09的压缩率相较于 GI（使用量化顶点） 取得领先。因此，laced wires 是唯一能够减少每个顶点存储大小 的方法，相较于紧凑型 GI 格式更加高效。如果增加额外的顶点属性，并利用网格单元的局部坐标系，则可以进一步摊销网格单元额外 9 字节的存储开销。与连接关系数据处理方式类似，$n_{v}/(n_{t} + n_{v})$的通用数据（无法明确归类为连接关系或几何数据）被归入几何数据部分。

TS 排名最末，因为它的冗余顶点导致了 较大的额外开销。IX 和 SC 无法摊销附加数据（本地索引和边界顶点副本），导致它们的压缩率分别为 0.48 和 0.66。与 GI 每个顶点 5.63 字节 的存储需求相比，IX 需要 11.74 字节，SC 需要 8.51 字节，而我们的方法仅需 5.14 字节。

图 7 展示了各方法在空间/时间维度上的对比。Laced wires 占用的内存最少。IX 和 SC 解码和渲染速度更快，但需要显著更多的内存。TS 不具备可行性，因为它的 时间和内存开销过大。GI 作为标准格式，在空间和时间上均被 SC 和 laced wires 超越，表明压缩技术带来了明显的优势。

\begin{figure}[h]
    \centering
    \begin{minipage}{\linewidth}
        \centering
        \includegraphics[width=\linewidth]{11.png}
    \end{minipage}

    \vspace{0.5em}
    
    \begin{minipage}{\linewidth}
        \fontsize{11pt}{13.5pt}\selectfont
        \bfseries
        图7\quad 所有比较的网格和网格单元格式在空间-时间维度上的对比图。基本版本用方形表示，紧凑版本用圆形表示，我们的方法用三角形表示。
    \end{minipage}
\end{figure}

\subsection{顶点量化分析}

我们的场景包含的网格已经事先按最低精度进行了量化，以确保一致的视觉质量。为了展示这种事先量化对压缩的影响，我们将蜜蜂网格在全球范围内使用不同的每个顶点坐标位数进行量化，确定最大量化误差，并运行我们的算法对每个网格单元进行局部量化，确保所有结果都保持在最大全局量化误差之下。这个实验的结果在图 8 中展示，位数范围为 8 到 16 位每个顶点坐标。

\begin{figure}[h]
    \centering
    \begin{minipage}{\linewidth}
        \centering
        \includegraphics[width=\linewidth]{12.png}
    \end{minipage}

    \vspace{0.5em}
    
    \begin{minipage}{\linewidth}
        \fontsize{11pt}{13.5pt}\selectfont
        \bfseries
        图8\quad 在我们的算法中，当匹配全局量化的最大全局解码误差时，顶点大小的比较（以 B/v 为单位）。
    \end{minipage}
\end{figure}

有趣的是，当解码误差相匹配时，我们的局部量化和全局量化之间的顶点精度差异并不是恒定的。这有两个主要原因：（1）角点是在全局坐标系中量化的，即它们的量化位通常与全局量化匹配，导致平均顶点大小出现异常值；（2）解码误差并不是随着量化位数的增加而线性减少。尽管如此，对于给定的容差，所有结果都显著低于全局量化的顶点大小。

\subsection{与现有方法的对比}

在本节中，我们基于原始论文中报告的数据，与现有方法进行对比。由于无法获取这些方法的代码，为了使性能数据具有可比性，我们基于公开的 GPU 基准测试 [WMMS23] 估算了一个速度修正因子。我们关注的是连接信息的压缩率（以每个三角形的比特数 bpt 计算）以及解码和渲染的总时间（以 Giga-triangles per second, Gtps 计算，使用基础渲染方法，如 Phong 着色）。

我们将这些方法分为三类：第一类（Class 1）包括在 CPU 上进行解码的方法 [Goo22, GLLR13]，它们主要作为连接信息压缩率的参考（但不考虑速度）；第二类（Class 2）包括在 GPU 上解码到内存的方法 [MKSS12, JBG17]；第三类（Class 3）包括 GPU 上实时解码的方法，除我们的方法外，我们仅知道另外两种方法 [NFR22, KSW21]。

表 6 总结了这些数据。从中可以看出，第一类和第二类方法的连接信息压缩率接近理论最优值（1–2 bpt），但由于需要将解码后的数据存储到内存中，最大速度受到了限制。此外，这些方法依赖于解压后的网格能够完全装入内存。相比之下，第三类方法的连接信息占用更多的存储空间（12–17 bpt），但解码速度明显更快。

尽管本次对比中的性能数据是基于 GPU 速度的粗略估算，因此结果可能存在一定误差，但整体趋势清楚地表明，laced wires 方法在速度方面优于所有现有方法。这一观察结果可以归因于网格着色器中对紧凑性和局部性的综合利用。

\begin{figure}[h]
    \begin{minipage}{\linewidth}
        \fontsize{11pt}{13.5pt}\selectfont
        \bfseries
        表6\quad 现有方法对比，按类别分组：Class 1（CPU 解码）、Class 2（解码至 GPU 内存）、Class 3（实时解码）。“Conn.” 列表示压缩后每个三角形的平均比特数（bpt）。“Final” 列报告了估算的 Gtps 速度（解码 + 渲染时间），并乘以根据 GPU 速度计算的修正因子（Corr[ection] factor）。对于 Meyer 等人 [MKSS12]，由于其结果存在较大波动，我们报告了“最小”和“最大”性能值，因此外推的结果可能不太可靠。对于 Jakob 等人 [JBG17]，由于未报告渲染速度，我们保守地假设其渲染时间为零。对于 Karis 等人 [KSW21]，由于未披露所用 GPU 型号，我们假设使用的是 RTX 3080（2020 年最新一代）。
    \end{minipage}

    \vspace{0.5em}

    \centering
    \begin{minipage}{\linewidth}
        \centering
        \includegraphics[width=\linewidth]{13.png}
    \end{minipage}
\end{figure}

\section{限制}

我们的方法在处理大规模网格时，编码时间可能会变得相当长，对于大型场景，编码可能需要数十分钟到数小时。编码时间较长的主要原因是，为了构造 laced wires，需要对最优顶点路径进行穷尽搜索。另一个计算开销较大的步骤是，在给定误差容限的情况下，搜索最小的顶点编码，这涉及多个坐标系之间的转换。编码时间可以通过启发式剪枝搜索以及优化编码器来显著减少。目前，我们的编码器仅运行在单线程 CPU 上，但尚未尝试多线程或 GPU 实现。

我们当前的表示形式在存储大小和解码速度之间进行了权衡。特别是，wire 头部占据了一定的内存空间，未来可以通过使用分层包围盒（hierarchical bounding boxes）等方法进一步优化。此外，似乎可以使用更大的网格单元，但目前尚不清楚网格着色器的缓存大小对网格单元规模的限制。

目前，laced wire 仅限于编码流形（manifold）几何体，因为每条 wire 中的顶点仅被访问一次，并且 laces 假定每条边最多相邻两个三角形。对于非流形（non-manifold）网格，可以在预处理阶段将其拆分为多个流形子网格进行处理。然而，如果需要生成过多的子网格，这种方法可能会导致次优结果。

\section{结论}

现代 3D 应用中不断增长的几何数据量带来了巨大的挑战。在本文中，我们提出了一种用于端到端渲染压缩网格的方法，该方法具有多个优势。资源通常以压缩格式存储，尽管压缩能够显著减少非易失性存储器中的数据大小，但 CPU 仍需在将数据提交至 GPU 进行渲染之前执行解压缩。因此，GPU 的内存占用和带宽需求并未因压缩而减少。在 GPU 端进行解压缩（如 Nanite 所采用的方法）可以减轻 CPU 的解压缩负担，使其能够执行其他任务。然而，现有的 GPU 解压缩方法会在 GPU 内存中存储解压缩后的数据，并同时保留压缩数据（内部管理压缩缓冲区），这实际上增加了 GPU 内存的总压力。

为应对这些挑战，我们提出了一种基于网格局部一致性的新型分解方法，即 laced wires。Laced wire存储了拓扑连接信息和几何数据，避免了顶点重复，并且可以独立进行压缩、存储和传输。小规模的 laced wire 集合组成了一个网格单元，可以在解码之前进行高效剔除。解压缩过程在网格着色器中实时进行，并直接将解码结果传输至光栅化器。

我们已证明，在 GPU 上按需实时解压几何数据以用于光栅化在当前硬件上是可行的。随着内存带宽和计算带宽之间的差距不断扩大，解压缩所带来的相对计算开销将会随着时间的推移而降低。如果硬件能够提供适当的增强（例如对量化属性的支持），解压缩吞吐量可能会在较低成本下大幅提升。此外，未来研究可以探索如何利用 laced wire 加速光线追踪，这可能需要结合网格单元表示方式与空间加速数据结构。 